% Chapter 8 — keep it running: orchestration, quality, operations.

\gmzoomin{operate}
\note{Chapter 8: the machine around all of it --- the dashed box on the map.}

\begin{frame}{Every month, a robot redoes only what changed}
  \centering
  \begin{tikzpicture}[x=1cm, y=1cm,
      st/.style={circle, minimum size=7.5mm, inner sep=0pt, line width=1.2pt, fill=gmSurface, font=\scriptsize},
      lb/.style={font=\fontsize{6}{7}\selectfont, text=gmInk, align=center, text width=1.55cm},
      go/.style={-{Stealth[length=1.6mm]}, line width=.8pt, draw=gmBase, shorten <=.6mm, shorten >=.6mm}]
    \foreach \i/\c/\ic/\t in {0/collect/\faSitemap/find pages, 1/collect/\faGlobe/fetch pages,
                              2/read/\faFile*/read pages, 3/link/\faLink/link names,
                              4/link/\faGavel/keep what matters, 5/events/\faStream/tell the story} {
      \node[st, draw=gm@col@\c, text=gm@col@\c] (a\i) at (\i*2.25,1.4) {\ic};
      \node[lb, below=1mm of a\i] {\t};}
    \foreach \i/\c/\ic/\t in {0/picture/\faImages/find templates, 1/picture/\faEye/read templates,
                              2/picture/\faImage/read the meme's image, 3/link/\faBook/Wikidata facts,
                              4/graph/\faProjectDiagram/build \& publish} {
      \node[st, draw=gm@col@\c, text=gm@col@\c] (b\i) at (\i*2.25+1.1,-1.3) {\ic};
      \node[lb, below=1mm of b\i] {\t};}
    \foreach \i [count=\j] in {0,...,4} \draw[go] (a\i) -- (a\j);
    \draw[go] (a5.south) |- ($(b0.north)+(0,.62)$) -- (b0.north);
    \foreach \i [count=\j] in {0,...,3} \draw[go] (b\i) -- (b\j);
    \node[font=\scriptsize\bfseries, text=gm@col@operate, anchor=east] at (a0.west) {1st of the month\enspace};
  \end{tikzpicture}
  \vskip7mm
  {\large Each step does \textbf{only what is new or changed} --- and can stop and restart safely.}
  \note{
  \begin{itemize}
    \item Layman version: every month a robot repeats the whole journey we just made with Doge ---
      but only for the pages that are new or changed. Doge, unchanged, is not read again.
    \item The chain: find pages, fetch them, read them, link names, keep what matters, tell the
      story, find templates, read them, read the meme's own image, fetch Wikidata facts, build and
      publish the graph. Each step starts the next.
    \item If anything fails half-way, running it again finishes the job without redoing or losing
      work. That is the property everything is built around.
  \end{itemize}}
\end{frame}

\begin{frame}{The machine}
  \centering
  \resizebox{!}{6.0cm}{\begin{tikzpicture}[x=1cm, y=1cm,
      bx/.style={rounded corners=1.6mm, draw=gmBase, line width=.9pt, fill=gmSurface, align=center, font=\scriptsize,
                 minimum height=.95cm, inner sep=2pt},
      go/.style={-{Stealth[length=1.6mm]}, line width=.7pt, draw=gmBase},
      grp/.style={rounded corners=2.5mm, draw=gmHair, line width=.8pt, fill=gmTint}]
    % groups
    \draw[grp] (-0.4,-0.7) rectangle (5.0,3.3);  \node[font=\scriptsize\bfseries, text=gmInk2, anchor=north west] at (-0.3,3.25) {Airflow 3 --- \FDags\ DAGs};
    \draw[grp] (5.6,-0.7) rectangle (10.0,3.3); \node[font=\scriptsize\bfseries, text=gmInk2, anchor=north west] at (5.7,3.25) {stores};
    % airflow
    \node[bx, draw=gm@col@operate, minimum width=2.1cm] (sch) at (1.05,2.2) {scheduler};
    \node[bx, draw=gm@col@operate, minimum width=2.1cm] (wk) at (3.5,2.2) {workers\\{\tiny Celery, Redis}};
    \node[bx, minimum width=2.1cm] (pg) at (1.05,0.4) {Postgres\\{\tiny run history}};
    \node[bx, minimum width=2.1cm] (dag) at (3.5,0.4) {stage code\\{\tiny library $\leftarrow$ store $\leftarrow$ DAG}};
    % stores
    \node[bx, draw=gm@col@read, minimum width=1.9cm] (mo) at (6.75,2.2) {MongoDB\\{\tiny one owner per collection}};
    \node[bx, draw=gm@col@graph, minimum width=1.9cm] (n4) at (8.95,2.2) {Neo4j};
    \node[bx, draw=gm@col@graph, minimum width=1.9cm] (fu) at (8.95,0.4) {Fuseki};
    \node[bx, minimum width=1.9cm] (fi) at (6.75,0.4) {files\\{\tiny builds, images}};
    % outside
    \node[bx, draw=gm@col@collect, minimum width=2.6cm] (web) at (2.3,4.5) {Know Your Meme, imgflip\\{\tiny via a rendering service}};
    \node[bx, draw=gm@col@events, minimum width=2.6cm] (gpu) at (7.8,4.5) {the lab's GPUs\\{\tiny language and vision models}};
    \node[bx, draw=gm@col@numbers, minimum width=4.6cm] (px) at (4.8,-1.9) {one door: port 8080\\{\tiny dashboard, Neo4j Browser, SPARQL, admin}};
    \draw[go] (web) -- (wk);  \draw[go] (wk) -- (gpu);
    \draw[go] (wk) -- (mo);   \draw[go] (wk.south east) -- (fi.north west);
    \draw[go] (mo) -- (n4);   \draw[go] (fi) -- (fu);
    \draw[go] (px.north) -- ++(0,.75);
  \end{tikzpicture}}
  \note{
  \begin{itemize}
    \item 15 services in Docker Compose on one server. Airflow 3 orchestrates 12 DAGs: a scheduler,
      Celery workers with Redis, Postgres for Airflow's own history.
    \item MongoDB holds every stage's output; each collection has exactly one owning module (the
      layering rule: pure library $\leftarrow$ store module $\leftarrow$ DAG; a DAG never issues a
      query, a library never touches a database).
    \item The graph goes to Neo4j and Fuseki, plus files. Models run on the lab's GPU servers
      (Ollama behind Open WebUI): ministral-3:14b for text, qwen3-vl:32b for images.
    \item One door: the campus firewall lets only port 8080 through, so a Caddy proxy serves the
      Streamlit dashboard, Neo4j Browser, the SPARQL endpoint and the admin tools behind it.
  \end{itemize}}
\end{frame}

\begin{frame}{Principles that made it hold}
  \begin{columns}[t]
    \column{.5\textwidth}
      \begin{itemize}\setlength{\itemsep}{2.6mm}
        \item \textbf{Stamps, not timestamps}: a stage redoes work when the code, prompt, model,
          lexicon or input that produced it changed
        \item \textbf{Dead letters}: what fails is kept aside with its reason, never retried blindly
        \item \textbf{Models on record}: every output names the model that produced it
        \item \textbf{Every run leaves a summary}: the dashboard reads them
      \end{itemize}
    \column{.44\textwidth}
      \begin{tikzpicture}[baseline=(current bounding box.north)]
        \node[fill=gmTint, rounded corners=2mm, inner sep=3.5mm, text width=5.2cm, align=left] {
          \gmwhy{Shared GPUs}\par\vskip1.5mm\small
          The models run on the lab's servers, one call at a time. Reading 36,673 stories and
          26,868 pictures took days --- so every step had to survive interruption and
          resume where it stopped.};
      \end{tikzpicture}
  \end{columns}
  \note{
  \begin{itemize}
    \item Staleness by stamps: each output records the versions that produced it (parser version,
      prompt version, schema hash, model, lexicon version, linker version\dots). Selection is
      ``missing or stale under the current stamps''. Changing the event prompt re-extracts events;
      it does not re-scrape pages.
    \item Dead letters: parse failures, event failures, template failures are stored with their
      reason, and not retried until something relevant changes.
    \item Models on record: every output stores the model that produced it, and the curation judge
      never falls back to another model (event extraction still may, by design --- recorded either
      way), so the graph never mixes readings without saying so.
    \item Run summaries in Mongo feed the dashboard: progress, coverage, failures, per stage.
    \item The GPU constraint shaped everything: backfills of days, one call at a time.
  \end{itemize}}
\end{frame}

\begin{frame}{Checked, measured, written down}
  \begin{columns}[t]
    \column{.31\textwidth}
      \gmbig{\FTests}{automated tests}\par\vskip1.5mm
      {\scriptsize\color{gmInk2}\FTestLines\ of the \FPyLines\ lines of Python; saved pages and model answers as fixtures}
    \column{.31\textwidth}
      \gmbig{\FGaps}{technical gaps logged}\par\vskip1.5mm
      {\scriptsize\color{gmInk2}what is weak, why, and what would fix it --- in the repository, in the open}
    \column{.31\textwidth}
      \gmbig{\FCommits}{commits since \FStart}\par\vskip1.5mm
      {\scriptsize\color{gmInk2}from a scraper to a published, versioned graph}
  \end{columns}
  \vskip8mm
  {\small\color{gmInk2}Every model output is reviewed in a loop --- sample, read, fix, read again --- on
  contact sheets, blind audits and held-out sets never used for tuning.}
  \note{
  \begin{itemize}
    \item 1,195 tests, about 14k of the 44k lines of Python. Parsers are tested on saved pages;
      model-facing code on recorded answers; the DAG chain itself is tested (who triggers whom,
      with which parameters).
    \item 14 technical gaps logged in the repository: each says what is weak, what it costs, and
      the fix --- honest debt-tracking. Several chapters today closed one.
    \item Review loops: for every model layer, a sample is read (contact sheets of templates,
      blind re-reading of images, event review page), the code is fixed, and a holdout that was
      never used for tuning confirms the fix.
    \item 96 commits since 24 June: the whole project in about three and a half months.
  \end{itemize}}
\end{frame}

\gmzoomout{operate}
